% Modelo de Dissertação em Latex para o PPG em Engenharia Mecânica da UERJ
% Este modelo foi adaptado da versão disponibilizada no site da Engenharia Elétrica da UERJ
% http://www.pel.uerj.br/publico/Modelo_LaTeX_Dissertacao_UERJ.rar
% http://www.pel.uerj.br/defesas/
%
% Utilizei o WinEdt 6.0 com Miktex 2.9
%
% Para gerar o PDF usei a opção PDFTeXfy com o documento [masterthesis.tex] aberto e em foco.
%
% Não consegui usar as figuras em EPS como o modelo original. Usei PNG e JPG sem problemas.
%
% Felipe M. - 20/06/2012
%

% Modificações e adaptações ao modelo da UFPE feitas por João Cordeiro.
% Versão original disponibilizada como template do latex.
% Este modelo foi criado para servir de exemplo de formato, não devendo seu código ser
% utilizado como referência de formatação.
% PLAIN: O código é uma adaptação de uma adaptação. Funciona, mas está cheio de gambiarras.
%
% João Cordeiro - 05/03/2019
%



\documentclass[a4paper,12pt,oneside,openany,table,xcdraw]{uerj}

\usepackage[portuges]{babel}
%\usepackage[latin1]{inputenc}
\usepackage[utf8]{inputenc}
\usepackage{enumerate}
\usepackage{cite}
\usepackage{epsf,epsfig,psfig}
\usepackage{pagina}
\usepackage{indentfirst}
\usepackage{theorem}
\usepackage{fancyhdr}
\usepackage{setspace}
\usepackage{boxedminipage}
\usepackage{float}
\usepackage{makeidx}
\usepackage{amsmath}
\usepackage[hidelinks]{hyperref}
\usepackage{todonotes}
\usepackage{caption}
\usepackage{graphicx}
\usepackage{multirow}
\usepackage{blindtext}


\makeindex


\newtheorem{deff}{Definição}[section]
\numberwithin{equation}{chapter}

\theoremstyle{plain}

\bibliographystyle{abnt-num}

\newcommand*{\captionsource}[3]{%
    \caption{\label{#1}#2.}
    \vspace{-0.3cm}
    \caption*{Fonte: #3.}
}


\begin{document}

\hypersetup{
    colorlinks,
    citecolor=black,
    filecolor=black,
    linkcolor=black,
    urlcolor=black,
    linktoc=all
}
% Infos do projeto
%

\newcommand{\thecenter}{Centro de Tecnologia e Geociências da Universidade Federal de Pernambuco}
\newcommand{\thedepartment}{Departamento de Engenharia Elétrica}
\newcommand{\thecourse}{Curso de Engenharia de Controle e Automação}
\newcommand{\thetitle}{Título do Trabalho}
\newcommand{\thetype}{Trabalho de Conclusão de Curso de Graduação}
\newcommand{\theproftitle}{Bacharel em Engenharia de Controle e Automação}
\newcommand{\thestudent}{João Victor Cordeiro Coutinho}
\newcommand{\theadvisor}{Prof. Wolfgang Hoenig}
\newcommand{\thecity}{Recife}

% \thispagestyle{empty}\listoftodos\pagebreak
\thispagestyle{empty}\input{./01_Pre_textuais/Capa}
\pagebreak\thispagestyle{empty}\input{./01_Pre_textuais/FolhaDeRosto}
\pagebreak\thispagestyle{empty}\input{./01_Pre_textuais/Agradecimento}
\pagebreak\thispagestyle{empty}\input{./01_Pre_textuais/Epigrafe}
\pagebreak\thispagestyle{empty}\input{./01_Pre_textuais/Resumo}
\pagebreak\thispagestyle{empty}\input{./01_Pre_textuais/Abstract}

\fancypagestyle{plain}{
\fancyhf{} % clear all header and footer fields
\renewcommand{\headrulewidth}{0pt}
\renewcommand{\footrulewidth}{0pt}}
\pagestyle{plain}

\pagebreak

\def\listfigurename{LISTA DE FIGURAS}\listoffigures
\def\listtablename{LISTA DE TABELAS}\listoftables
\input{./01_Pre_textuais/ListaDeSiglas.tex}    % não coloquei LISTA DE SIGLAS no meu trabalho, mas fica aqui a chamada comentada.
\def\contentsname{SUMÁRIO}\tableofcontents

\fancypagestyle{plain}{
\fancyhf{} % clear all header and footer fields
\fancyhead[R]{\thepage}
\setlength{\voffset}{-1cm}
\setlength{\headsep}{1cm}
\renewcommand{\headrulewidth}{0pt}
\renewcommand{\footrulewidth}{0pt}}

\pagestyle{plain}

\pagebreak
\input{./02_Textuais/introducao/introducao.tex}

\input{02_Textuais/desenvolvimento/problemas.tex}

\input{./02_Textuais/desenvolvimento/solucao_proposta.tex}

\input{./02_Textuais/desenvolvimento/tests/clinical_test.tex}


% inserir demais capítulos aqui
% -----------------------------
% -----------------------------
% -----------------------------
% -----------------------------





\pagebreak
\input{./02_Textuais/conclusao/conclusao.tex}

\pagebreak
\def\bibname{REFERÊNCIAS}




% abaixo segue a chamada para o arquivo [.BIB]. Utilizei o programa JABREF para montar o arquivo com minhas referências.
% \bibliography{dissertacao}
\bibliography{tcc}





%felipe% \printindex    %Removi o índice remissivo para a versão oficial do trabalho.


\end{document}